\section{架构执行与品味守护}

\subsection{机械化的架构约束}

文档本身无法让一个完全由 Agent 生成的代码库保持一致性。团队通过\textbf{执行不变量（invariants），而非微管理实现}的方式来解决这一问题。

\begin{importantbox}{严格的分层架构模型}
每个业务域被划分为固定的层次集合，层间依赖方向严格验证，只有有限的合法边集。具体分层为：

\textbf{Types $\rightarrow$ Config $\rightarrow$ Repo $\rightarrow$ Service $\rightarrow$ Runtime $\rightarrow$ UI}

跨切面关注点（auth、connectors、telemetry、feature flags）通过单一显式接口 \textbf{Providers} 注入。其他任何依赖关系都被禁止，并通过机械化手段强制执行。
\end{importantbox}

文章特别指出：\textbf{这种严格的架构通常是拥有数百名工程师的大型团队才会考虑的事情}。但在 Agent 开发中，它是早期前提——正是这些约束使得高速开发不会导致衰退或架构漂移。

\subsection{自定义 Linter 与品味不变量}

团队使用自定义 Linter（同样由 Codex 生成）和结构化测试来执行规则：

\begin{table}[H]
\centering
\small
\begin{tabular}{p{4.5cm}p{8.5cm}}
\toprule
\textbf{强制规则类型} & \textbf{示例} \\
\midrule
结构化日志 & 静态强制所有日志输出遵循结构化格式 \\
命名约定 & Schema 和类型的命名必须遵循统一规范 \\
文件大小限制 & 防止单文件过大，保持可读性 \\
平台特定可靠性要求 & 针对不同平台的性能和可靠性基线 \\
数据边界验证 & 要求在数据边界解析数据形状（如使用 Zod） \\
依赖方向验证 & 代码只能``向前''依赖固定的层次顺序 \\
\bottomrule
\end{tabular}
\caption{品味不变量（Taste Invariants）示例}
\end{table}

\begin{warningbox}{自定义 Linter 的关键设计：错误消息即 Agent 指令}
因为 Linter 是自定义的，团队将错误消息设计为包含\textbf{修复指导}（remediation instructions），这些指导会直接注入 Agent 的上下文中。这意味着 Lint 错误不仅告诉 Agent ``你做错了什么''，还告诉它``应该怎么修''。在人类开发中，这些规则可能显得教条；在 Agent 开发中，它们是乘数效应——一旦编码，即可在所有地方同时生效。
\end{warningbox}

\subsection{约束与自由的平衡}

团队的理念类似于领导一个大型工程平台组织：

\begin{itemize}[nosep]
  \item \textbf{中央强制执行}：边界、正确性、可复现性
  \item \textbf{局部允许自治}：在约束范围内，Agent 有充分自由决定解决方案的表达方式
\end{itemize}

结果是代码并不总是符合人类的风格偏好——\textbf{但这没关系}。只要输出是正确的、可维护的、且对未来的 Agent 运行可读，就达到了标准。

人类的品味通过以下渠道持续反馈到系统中：

\begin{enumerate}[nosep]
  \item Review 评论
  \item 重构 Pull Request
  \item 面向用户的 Bug 被捕获为文档更新
  \item 当文档不够时，将规则提升为代码（Linter）
\end{enumerate}

\subsection{本章小结}

在 Agent-First 开发中，架构约束不是事后添加的技术债务治理工具，而是\textbf{速度的前提条件}。通过严格的分层架构、自定义 Linter（带修复指导的错误消息）、以及``中央强制边界、局部允许自治''的平衡策略，团队在保持高速产出的同时防止了架构衰退。

\newpage

